%% kapa: exam=1 round=34 year=2023 period=2 form=A subject=accounting subject_ko=회계학 seq=3 number=43
%% source: Q-Net 감정평가사 기출문제 (cst00309) artlSeq=5211789 file=raw/1008/5211789/2023년도 제34회 감정평가사 1차시험 2교시.pdf
%% conversion: kordoc 4.24.51 -> tools/gen1.py (auto); review=proofread-sonnet; flags=none
\begin{question}{43}
투자부동산의 분류에 관한 설명으로 옳지 않은 것은?
\begin{choices}
  \item 미사용부동산을 운용리스로 제공한 경우에는 투자부동산으로 분류한다.
  \item 리스계약에 따라 이전받은 부동산을 다시 제3자에게 임대한다면 리스이용자는 해당 사용권자산을 투자부동산으로 분류한다.
  \item 지배기업이 다른 종속기업에게 자가사용 건물을 리스하는 경우 당해 건물은 연결재무제표에 투자부동산으로 분류할 수 없다.
  \item 건물 소유자가 그 건물의 사용자에게 제공하는 부수적 용역의 비중이 경미하면 해당 건물을 투자부동산으로 분류한다.
  \item 처분예정인 자가사용부동산은 투자부동산으로 분류한다.
\end{choices}
\end{question}
\answer{5}
